.bp41
.he
.mt7
.mb9
.po8
.ll65
.ig
CHAPTER 9
T H E	 A C C O U N T	  F I L E
A N D	 G L	I N T E R F A C E
.sp2
.pp5
At first time startup the Payroll System moves directly to
the Account Entry program after Parameter Entry is complete.
If you have chosen to use the default Payroll System (as explained in
Chapter 7), no requests are issued and the default account file is
created automatically.	If you have chosen to do further customization,
the Account Entry screen appears so you can add or change account
information.
Once the startup sequence
is complete and the account file created,
you enter or change account
information by choosing the Account Entry program from the system menu.
.he Reference: Chapter 9
.pp
The default Payroll Account File contains four accounts (Figure 5.1).
.sp
.cp9
.li
      Payroll System				 (General Ledger)
 Account Reference Number     Account Name	  Account Number

	   1		 CASH			      011100
	   2		 ACCRUED LIABILITY	      121100
	   3		 SALARY EXPENSE 	      411100
	   4		 ACCRUED PAYROLL COSTS LIAB.  131100

	      (Figure 5.1:  The Default Account File)
.sp2
These accounts are the bare minimum required for payroll processing.
Account number 1 is the cash account credited when checks are written.
If you do not select the automatic check writing feature you may want
to change this to an accounts payable account.	Account number 2 is the
liability account credited when you incur a liability by withholding money
from an employee's wages for taxes or other deductions.  Account number
3 is the salary expense account debited for the total cost of the
employee to the company (employee gross plus employer payroll costs).
If you choose the option to distribute employee cost to more than one
account, additional accounts (e.g., Admin. Salaries, Programming Costs,
Collins Project--Salaries, etc.) would be used for distributing
the cost of the employee to the employer.   The fourth account, Accrued Payroll
Costs Liability, is debited for the employers
payroll expenses, such as Company
FICA and Company FUTA.
.pp
Whether you choose GL interface or not, all account numbers must conform
to the following numbering scheme:
.sp
.in5
.ti-3
1. The first digit of a CASH account must be 0 (zero), e.g. 011100.
.sp
.ti-3
2. The first digit of a LIABILITY account must be 1 (one), e.g., 121100.
.sp
.ti-3
3. The first digit of an EXPENSE account must be 4 (four), e.g., 411100.
.qi
.sp
.pp
The four default accounts, plus additional accounts
you may enter should always be
referred to by their Payroll System account reference numbers.	The Payroll
System account reference number is the relative position of the account
record on the Account file.  For example, the first
account on the account file has account reference number 1,
the second account has account reference number 2, and so on.
.pp
If you select
the General Ledger interface option you must enter a valid four or
six digit General Ledger account number for each account. The fourth digit
from the left may never be zero, since that would mean it was a title
account and not a balance holding account (see the GL Reference Manual).
The General
Ledger data disk with the chart of accounts file must be mounted on one
of the disk drives when you enter the GL account numbers so the Accont Entry
program can verify the accuracy of the account information you key in.
The Account Entry program prompts you to insert the GL data disk at
the proper time.
.pp
If you do not
selected the General Ledger interface option you may enter any number of up to
six digits, subject to the first digit numbering restrictions explained above.
A printout showing the account numbers and names currently
in use can be generated from the system menu at any time.
.pp
Systemwide and employee specific deductions, and company payroll costs
such as CO FUTA and CO SUI should be liability accounts (begin with 1).
The accounts used to distribute the employees cost to the company should
be expense accounts (begin with 4).  If the accounts are numbered
improperly an imbalance will show on the Payroll Journal Ledger Account
Summary.  The GL batch will also be incorrect.
.pp
Before you can respond to any Payroll System request for an account number
with an account
number other than one already on the account file, the new account name and
number
must first have been entered onto the file through the Account Entry program.
This is because
the Payroll System checks the account information you enter against
the account file before it accepts the entry
to make sure you do not enter an invalid or non-existent
account.
.he
.bp
.ig
CHAPTER 10
S Y S T E M - W I D E
E A R N I N G S    A N D   D E D U C T I O N S
.sp2
.pp5
At first time startup, after Parameter Entry and Account Entry, the
Payroll System moves directly to the program used to enter system-wide
earnings and deduction information.  If you have choosen to use the
default Payroll System (as explained in Chapter 7), the System Deduction
Entry program issues no requests and creates the default System Deduction
and Earnings file automatically (see Figures 10.1, 10.2, and 10.3 for the
values in the default file).
Once the startup sequence is complete and the System Deduction and Earnings
file created,
you call this program directly from the menu.
The System Deduction Entry Program
is used for entering standard withholding and deduction information, changing
the Federal tax table, and creating up to five
system-wide earnings or deductions of your choice.
.he Reference: Chapter 10
.sp
.ce
STANDARD DEDUCTION RATES
.pp
The System Deduction Entry program requests the following
information:
.sp
.in5
.ti-5
FWT: Whether Federal income tax is to be withheld from employee paychecks,
and if so, what account should be credited with the accrued liability.
.sp
.ti-5
SWT: Whether state income tax is to be withheld from employee paychecks,
and if so, what account should be credited with the accrued liability.
.sp
.ti-5
LWT: Whether local income tax is to be withheld from employee paychecks,
and if so, what account
should be credited with the accrued liability.
.sp
.ti-5
FICA: Whether Social Security tax is
to be deducted from an employee's earnings,
the account to be credited with the accrued liability,
the rate at which it is to be deducted, and the highest total
salary subject
to the deduction.
.sp
.ti-5
CO FICA: Whether the company must pay into the Social Security fund for
an employee, the account to be credited with the accrued liability,
the rate of the
employer's contribution, and the highest total salary subject to the tax.
.sp
.ti-5
SDI: Whether a deduction should be made from an employee's gross pay
for State Disability Insurance, the account to be credited with the
accrued liability, the rate at which the deduction is to be made,
and the highest total salary subject to the deduction.
.sp
.ti-5
CO FUTA: Whether the company must pay into the Federal Unemployment
Tax fund, the account number to be credited with the accrued
liability, the rate
at which the tax is calculated, and the highest total
salary subject to the tax.
.sp
.ti-5
FUTA MAX STATE CREDIT: The highest rate that can be credited against the
CO FUTA rate if the employer makes a contribution to a state unemployment
insurance fund.  For example, if the CO FUTA rate is 3.4, the CO SUI
rate 3.0, and the largest rate that can be credited against the
FUTA liability is 2.7 (FUTA MAX STATE CREDIT), the effective FUTA
rate will be 0.4 percent (3.4 - 2.7).  If the CO FUTA rate is 3.4,
the CO SUI rate 2.0, and the FUTA MAX STATE CREDIT is 2.7, the effective
FUTA rate will be 1.4 percent (3.4 - 2.0).
.sp
.ti-5
CO SUI: Whether the company must pay into a state unemployment tax fund,
the account to be credited with the accrued liability, the rate at which the
contribution is calculated, and the highest total salary subject to the tax.
.sp
.ti-5
EIC: Whether the company must make an advance payment for the earned
income credit to employees who request it,
the liability account to be offset by the advance payment,
the rate at which the payment is calculated, and the highest annual equivalent
wage to which
that rate applies.
.sp
.ti-5
EIC EXCESS: Provision has been made for entering the rate and limit
used for calculating the reduction in the earned income credit payment
that results when an employee's annual equivalent wage exceeds the
EIC EXCESS limit.
.qi
.cp18
.sp
.po12
.li
		 STANDARD DEDUCTION RATES

		  ACCT
	   USED?   NO	      RATE	    LIMIT

   FWT:     [Y]   [  2]
   SWT:     [Y]   [  2]
   LWT:     [N]   [  2]
   FICA:    [Y]   [  2]    [   6.13]  [   22,900.00]
   CO FICA: [Y]   [  4]    [   6.13]  [   22,900.00]
   SDI:     [N]   [  2]    [   1.00]  [    6,000.00]
   CO FUTA: [Y]   [  4]    [   3.40]  [    6,000.00]
   FUTA MAX STATE CREDIT:  [   2.70]
   CO SUI:  [Y]   [  4]    [   0.00]  [    6,000.00]
   EIC:     [N]   [  2]    [  10.00]  [    5,000.00]
   EIC EXCESS:		   [  12.50]  [    6,000.00]

    (Figure 10.1: Default Standard Deduction Rates)
.po8
.sp2
.ce
FEDERAL WITHHOLDING TAX TABLE ENTRY
.pp
Applewood Computers distributes the Payroll System with the current
Federal withholding tax table. Income tax withholding is based on the
percentage method (see IRS circular E, Employer's Tax Guide), using the
ANNUAL payroll period table.  The System Deduction Entry program lets you
to change this table when necessary.  See Chapter 22 for details.
.sp
.cp13
.po13
.li
	      FEDERAL WITHHOLDING TAX TABLE ENTRY

	       ALLOWANCE AMOUNT:  [    1,000.00]

	   M A R R I E D		S I N G L E
	CUTOFF	    PERCENT	    CUTOFF	 PERCENT
1. [   2,400.00] [     15.00]	[   1,420.00] [     15.00]
2. [   6,600.00] [     18.00]	[   3,300.00] [     18.00]
3. [  10,900.00] [     21.00]	[   6,800.00] [     21.00]
4. [  15,000.00] [     24.00]	[  10.200.00] [     26.00]
5. [  19,200.00] [     28.00]	[  14,200.00] [     30.00]
6. [  23,600.00] [     32.00]	[  17,200.00] [     34.00]
7. [  28,900.00] [     37.00]	[  22,500.00] [     39.00]

	   (Figure 10.2: Default Federal Tax Table)
.sp2
.po8
.ce
SYSTEM EARNING AND DEDUCTION INFORMATION
.pp
The Payroll System lets you define up to
five system-wide earnings or deductions categories
of your choice.  You enter a description, state whether it is
an earning or deduction, and state in which box on the paycheck stub it is
to be printed.
If a deduction, you also enter the account number to be
credited with the
liability. You state whether the earning or deduction is a percentage
of the employee's salary or a flat amount each pay period.
You can also state whether there is a limit to the amount that can
be earned or deducted in one year, and if so, enter that amount.
If the entry is an earning, you can assign a tax exclusion code to it that
indicates one of the following:
.sp
.in8
.ti-3
1. Subject to all taxes and percentage based deductions.
.ti-3
2. Subject to all taxes and percentage based deductions except FICA.
.ti-3
3. Subject to all taxes and percentage based deductions
except Federal, state, and local income taxes.
.ti-3
4. Not subject to taxation or percentage based deductions.
.qi
.sp2
.po4
.li
	       SYSTEM EARNING AND DEDUCTION INFORMATION

   USED?    DESC       E/D ACCT A/P  FACTOR   LIMITED  LIMIT   XLCD CAT

1. [N] [	     ] [E] [  ] [A] [	  0.00] [Y] [	  0.00] [1] [2]
2. [N] [	     ] [E] [  ] [A] [	  0.00] [Y] [	  0.00] [1] [2]
3. [N] [	     ] [E] [  ] [A] [	  0.00] [Y] [	  0.00] [1] [2]
4. [N] [	     ] [E] [  ] [A] [	  0.00] [Y] [	  0.00] [1] [2]
5. [N] [	     ] [E] [  ] [A] [	  0.00] [Y] [	  0.00] [1] [2]

    (Figure 10.3: Default System Earning and Deduction Information)
.br
.po8
.he
.bp
.ig
CHAPTER 11
P U T T I N G	E M P L O Y E E S   O N   T H E   S Y S T E M
.sp2
.pp5
Putting employees on the Payroll System is a three step procedure:
.he Reference: Chapter 11
.sp
.in8
.ti-3
1. Enter the employee's name, address, social security number,
and other information through the Employee Entry Program.
.sp
.ti-3
2. Enter earnings, deduction, exemption, and costing information for the
employee through the Employee Ded/Earn and Cost program.
.sp
.ti-3
3. Enter pay interval, pay rate, withholding allowance,
vacation, sick leave, and comp time information
for the employee through the Employee Rate Entry program.
.qi
.sp
Step one must
be completed for an employee before step two or
step three can be taken for
that employee.	This is simply to say that before pay rate or
deduction information for an employee can be entered, a record for
the employee must be created by the Employee Entry program.
.sp2
.ce
EMPLOYEE ENTRY
.pp
When you select the Employee Entry program from the system menu,
an entry screen appears with the cursor waiting
at the Employee Number field.  From there you can give commands
to add new employees
to the file, or to change or examine information for employees
already on the file.
.pp
The Employee Entry program requests the following information:
.in5
.sp
.ti-5
EMPLOYEE NUMBER: Unless you are entering information for an employee
for the first time, in which case the program assigns the next available
number to the employee automatically,
you may enter the five digit number of an employee
whose record you wish to examine or change.  To examine
or change information for an employee, instead of accessing
the employee's record by the employee number, you can
issue a series of
Next-Record or Previous-Record commands to "thumb" through the file.
.sp
.ti-5
NAME:
.ti-5
ADDRESS: Two lines are available for the street address.
.ti-5
CITY:
.ti-5
STATE:
.ti-5
ZIP:
.ti-5
PHONE:
.ti-5
SOCIAL SECURITY NUMBER:
.ti-5
BANK ACCOUNT NUMBER: This is an optional entry for informational
purposes only.
.ti-5
BIRTH DATE:
.ti-5
SEX:
.sp
.ti-5
PENSION (Y/N): State whether the employee is covered under a pension
plan.  This information is required for the W-2 form.
.sp
.ti-5
JOB CODE: This is an optional entry for informational purposes only.
.sp
.ti-5
START DATE: The date the employee started working for the company.
.sp
.ti-5
TERM. DATE: If an employee left or was terminated, enter the date he or she
left the company's employ.
.sp
.ti-5
EMPLOYMENT STATUS: You can designate an employee as Active, On Leave, or
Terminated. Automatic pay generation is suppressed for employees on
leave.
The year-end closing program deletes
terminated employees from the file after a W-2 form has been printed, and
makes the number available for reassignment. You can print W-2 forms for
terminated employees at any time.
.qi
.sp2
.ce
EMPLOYEE DEDUCTION/EARNINGS AND COST DISTRIBUTION
.pp
When you select the Employee Deductions/Earnings and Cost Dist.
program from the system
menu, an entry screen appears with the cursor
waiting at the Employee Number field.  You can then enter an employee
number, or issue a Next-Record or Previous-Record command until
you find the desired employee record.  After the program displays the
employee's name and social
security number, you can make entries or changes to the
following requests:
.in5
.sp
.ti-5
GL COST DISTRIBUTION OF EMPLOYEE GROSS AND EMPLOYER COST: The sum of
all employee pay (including tips collected but excluding tips reported
and EIC payments)
and additional payroll expenses paid by the employer for
the employee can be distributed to up to five accounts if
the option is requested in Parameter Entry.  If you do not want a
particular pay type to be so distributed, you may have to make an
adjusting entry to your General Ledger. To set up the distribution accounts
you enter the Payroll System Account Reference
number (see Chapter 9) and the
percentage of the total cost of the employee to the employer that is
to be allocated to that account.  Thus, if a particular employee works
part of the time in one department and part in another, the cost
of that employee can be distributed to the appropriate department
or project for accounting purposes.
.sp
.ti-5
EXEMPTIONS FROM TAX WITHHOLDING:  A series of YES or NO requests
lets you exempt a specific employee from Federal, state,
local, FICA, or SDI withholding.
.sp
.ti-5
EXEMPTIONS FROM SYSTEM DEDUCTIONS:  Five YES or NO requests
permit you to exempt specific employees from any of the
five system-wide earnings or
deductions you can create through the
System Deduction Entry program (see Chapter 10).
.sp
.ti-5
EMPLOYEE SPECIFIC DEDUCTIONS:  These Payroll System requests let
you define up to three earnings or deductions for
a specific employee.  You enter a description, state whether it is
an earning or deduction, and state in which box on the paycheck stub it
is to be printed.  If a deduction, you enter the account number to be
credited with the liability. If an earning, the expense is debited
to the account or accounts specified for distributing the employee's
cost to the company.
You can state whether the earning or deduction is a percentage of the
employee's salary, or a flat amount each pay period.
You can also state if there is a limit to the amount that may be earned or
deducted in one year, and if so, enter the amount.
If the entry is an earning, you can state whether the
amount is:
.sp
.in12
.ti-3
1. Subject to all taxes and percentage based deductions.
.sp
.ti-3
2. Subject to all taxes and percentage based deductions
except FICA.
.sp
.ti-3
3. Subject to all taxes and percentage based deductions
except Federal, state, and local income taxes.
.sp
.ti-3
4. Not subject to taxation or percentage based deductions.
.qi
.sp2
.ce
EMPLOYEE RATE ENTRY
.pp
When you select the Employee Rate Entry program from the system menu,
an entry screen appears with the cursor waiting at
the Employee Number field.  You can then enter an employee number, or
issue Next-Record and Previous-Record commands until you find the
desired record. After the program displays
the employee's name and Social Security number,
you can enter or change the following information:
.sp
.li
PAY INTERVALS AND PAY RATES:
.in5
.sp
.ti-5
PAY TYPE:  State whether the employee is paid
by the hour or is a salaried employee.	This information affects only
certain report printing procedures.
.sp
.ti-5
PAY INTERVAL: State whether the employee is paid weekly, bi-weekly,
monthly, or semi-monthly.
.sp
.ti-5
PAY RATE (1-4): The Payroll System allows four types of pay rates.
The descriptive name given to each pay rate in
Parameter Entry appears beside
the current rate amount.  Usually these are "Regular Pay", "Overtime",
"Special Overtime", and "Other Pay", "Commission Pay", or "Piecework Pay",
but they can take any description.
.pp
Unless you enter an overriding Rate1 transaction through the Pay Transaction
Entry program, an employee is paid the Rate1 amount times the
number of Autogen Units established for that employee
(for example, $5.50 x 40 hrs, or $1,500 x 1 pay unit). If no Autogen Units
have been established for an employee, you must enter a RATE1 transaction
through the Pay Transaction Entry program for the number of units the
employee worked
if the employee is to receive any Regular (Rate1) pay.
.pp
A special RATE0 (Rate Zero) transaction can be made through the Pay
Transaction Entry program, so that
any pay units over the number of Normal Units established for that employee
are paid at Rate2 and all Normal Units are paid at RATE1.
The Rate2 and Rate3 default values
are Rate1 times the Rate2 and Rate3 Factors (entered through
Parameter Entry), respectively.
Of course, RATE2 and
RATE3 can be changed to any value you desire for an employee.
.sp
.ti-5
RATE4 TAX EXCLUSION TYPE: You can set the tax exclusion status for
individual employees according to the table below:
.sp
.qi
.cp10
.li
	    1. Subject to all taxes and
	       percentage based deductions.
	    2. Subject to all taxes and
	       percentage based deductions
	       except FICA.
	    3. Subject to all taxes and
	       percentage based deductions
	       except FWT. SWT, LWT.
	    4. Not subject to taxes or
	       percent deductions.
.in5
.sp
.ti-5
AUTOGEN UNITS: Many employees, especially those on a salary, receive
the same pay amount each pay period. Instead of entering a Rate1 transaction
each time for these employees, you can set the number of Rate1 units
for which pay will be automatically generated.	Of course, other types of
transactions can be entered for the employee (a bonus, advance on wages,
vacation taken, etc.), but if a Rate1 transaction is entered, automatic pay
generation will be suppressed and the employee paid for the number of Rate1
units entered.	You would normally set the number of
Autogen Units for a salaried employee to 1 (with the Rate1 rate set to
the amount the employee earns each pay period, such as $1,200.00).
.sp
.ti-5
NORMAL UNITS: This is the number of units an employee normally works
in one pay period.  If you enter a RATE0 transaction through the
Pay Transaction Entry program, the employee is paid at RATE2 for units
over the Normal Units figure.  Normal Units are usually expressed in hours
for hourly employees, pay periods for salaried employees, pieces for
piecework employees, etc.
.sp
.ti-5
MAXIMUM PAY: This is the most gross pay an employee may receive in one pay
period without an warning message being issued when the employee is processed
by the Apply Payroll and the Print Payroll Journal program.
The default value for this field comes from
multiplying the Rate1 amount times the Maximum Pay Factor entered
in Parameter Entry.
.qi
.sp
.li
TAX WITHHOLDING ALLOWANCES:
.in5
.sp
.ti-5
MARITAL STATUS: Married or Single. Needed for figuring
tax withholding.
.sp2
.qi
.li
FEDERAL WH ALLOWANCES:
STATE	WH ALLOWANCES:
.in5
.ti-5
LOCAL	WH ALLOWANCES: Enter the number of withholding allowances for
Federal, state, and local income tax.
.qi
.sp2
.li
STATE OF CALIFORNIA SPECIAL INFORMATION:
.in5
.sp
.ti-5
HEAD OF HOUSEHOLD: State whether the employee is considered a head
of household for California tax purposes.  This and the following request
are asked only if California (CA) is the state given in the company address
in Parameter Entry.
.sp
.ti-5
SPECIAL WITHHOLDING ALLOWANCES:  Enter the number of special withholding
allowances to which the employee is entitled under California's tax laws.
.sp2
.ti-5
VACATION RATE: Enter the number of days or hours of vacation the
employee accumulates each pay period.  To avoid confusion, enter this,
as well as sick leave and compensatory time off information, consistently
in terms of days or hours for all employees on the system.  If both hours
or days are used on the same payroll system,
an employee might accidently accumulate vacation time in hours, but use
it up in days or vice versa.
.sp2
.ti-5
VACATION ACCUMULATED: When first setting up your Payroll System,
enter the amount of vacation time the employee has acccumulated so far in the
year, or has remaining from the year before.
Once set,
this figure increases by the Vacation Rate each
pay period. Although it should not be necessary, you can change
the value in this field
at any time.
.sp
.ti-5
VACATION USED: This field keeps track of vacation used
during the current year.  Although you normally record Vacation Used
through the Pay Transaction Entry program, you can make changes to this
field directly if needed.
At the end of the year Vacation Used is subtracted from
Vacation Accumulated. The difference, if any, is placed in the
Vacation Accumulated field and the Vacation Used field is set to zero.
.sp
.ti-5
SICK LEAVE RATE: Works the same as Vacation Rate.
.sp
.ti-5
SICK LEAVE ACCUMULATED: Works the same as Vacation Accumulated.
.sp
.ti-5
SICK LEAVE USED: Works the same as Vacation Used.
.sp
.ti-5
COMP TIME ACCUMULATED: If an employee has earned compensatory time off during
the current year,
this field tells how much.  Comp time earned is recorded
by a special transaction
in the
Pay Transaction Entry program. At first time startup the
amount of comp time accumulated should be entered directly to this field.
Note that Comp Time Accumulated is not the same as comp time available
(remaining).  That figure results from subtracting Comp Time Used from
Comp Time Accumulated.
.sp
.ti-5
COMP TIME USED: This field tells how much compensatory time off an employee
has taken so far this year.  When an employee uses up comp time,
a comp time taken transaction should
be entered through the Pay Transaction Entry program, although the Employee
Rate Entry program
does permit direct access to this field.  At the end of the year
Comp Time Used is subtracted from Comp Time Accumulated.  The difference,
if any, is placed in the Comp Time Accumulated field, and the Comp Time
Used field set to zero.
.qi
.he
.bp
.ig
CHAPTER 12
T H E	P A Y R O L L	P R O C E S S I N G    C Y C L E
.sp2
.pp5
Once the Payroll System setup steps are complete,
you can begin the regular payroll
processing cycle. The payroll processing cycle is the sequence of
steps you repeat for every working period to compute earnings and deductions,
generate paychecks, and
update historical information.	For the purpose of explanation,
the quarter and year end processing steps are not considered part
of the payroll processing cycle, although they too are repeated
on a regular basis.  This chapter explains the payroll processing
cycle in general terms.  The following chapters cover steps in the cycle more
fully.
.he Reference: Chapter 12
.pp
In the simplest case (when all employees on the system are paid
on the same pay interval) the steps in the payroll processing cycle
are repeated once for every pay period.  Thus, if all employees are paid
weekly, you would repeat the following steps once a week:
.sp
.li
	     1. Enter Pay Transactions.
	     2. Print a Pay Transaction Audit Proof.
	     3. Back up all files.
	     4. Sort the Pay Transaction batch file.
	     5. Run the Apply Payroll program.
	     6. Print the Payroll Journal.
	     7. Print Paychecks and/or the Paycheck Register.
.pp
When weekly and bi-weekly employees are on the system, every other week
both are processed
by the same cycle.  Similarly, when semi-monthly and
monthly employees are on the system, both are processed by
the same cycle at the end of each month.  Note, however, that
although you may enter transactions for all four types into the same
transaction batch, the Apply Payroll program does not process
week-based and month-based employees in the same run.  Thus, the
Payroll Journal, paychecks, and Paycheck Register are always printed
separately for month-based and week-based employees, even if
they are to be paid on the same day.
.pp
When you run the Apply Payroll program, it asks whether
you want to process week-based (weekly and bi-weekly)
or month-based employees (monthly and semi-monthly).  If you
answer that you want to process week-based
employees, the Payroll Application program separates out the
month-based employee transactions and
stores them on the Pay Detail file.  Later, when you process
the month-based employees, the Payroll Application program will use
the transaction information stored on the Pay Detail file, as well
as any additional transactions for month-based employees on the current
transaction file.
.pp
If weekly and bi-weekly employees are on the system, when you instruct
the Payroll Application program to process week-based employees, it
will know whether to process only weekly employees, or both
weekly and bi-weekly employees by checking the week-base
"application number"
to see which were processed last time.	If only
weekly employees were processed during the last week-based application,
it knows to process weekly and bi-weekly employees next.  The same is true
for month-based employees.
.pp
When end of period processing should occur is determined not by how many
times you've repeated the cycle, but by the check date you enter at the
beginning of the Apply Payroll program.  This is the date that is
printed on each paycheck, and thus the date that determines which
quarter is to be charged for the payroll costs.  The date you enter is
checked against the Payroll System's internal calendar, and if end of
period processing is required, you are not allowed to run the
Payroll Application program until the end of period processing steps
are completed.
.sp2
.ce
ENTER PAY TRANSACTIONS
.pp
During this step you may enter any of a number of possible pay transactions for
an employee.  Among those possible are transactions for
Rate1 pay (Regular), Rate2 pay (Overtime), Rate3 pay
(Special Overtime), Rate4 pay (commission), vacation taken,
sick leave taken,
comp time taken or earned, bonus pay, tips collected, tips reported,
advance pay, advance repayment,
sick pay, vacation pay, expense reimbursement, or additional withholding.
Unless you enter a Rate1 pay transaction, an employee automatically
receives the Rate1 amount times the Autogen Units figure set for that employee.
If no automatic pay units have been established, the employee receives no
Rate1 pay.
.pp
Transactions may be entered for any employee, whether week-based or
month-based.  A transaction takes effect the next time the employee
is processed by the Apply Payroll program.  Any transactions not
used by the Apply Payroll program are stored on the Pay Detail file
until the next application for that interval.
.sp2
.ce
PRINT A PAY TRANSACTION AUDIT PROOF
.pp
Once the batch of pay transactions is complete and accurate, you must
generate a hardcopy transaction proof report as part of the
accounting audit trail.  Although a batch proof may be obtained as a working
proof
any time during the development of the batch, only the final proof is
considered the audit proof.
The audit proof must be taken before the file is sorted, and if
it is to retain it status as the audit proof, additional pay transactions
may not be entered.
.sp2
.ce
BACK UP ALL FILES
.pp
In order to process as many employees as possible,
the Apply Payroll program does not automatically create backup versions
of all the files it processes.
If the program is run at the wrong time, or if some difficulty
arises during or after processing, it may be very difficult to get back
to where you where before the application took place.  For this reason,
we strongly recommend that you make backup copies of all your disks before
going to the next processing step. Use the full-disk single density copy
program SDCOPY (see Chapter 35) provided with the
Payroll system if you are operating in
single density, use the full-disk copy program the manufacturer of your disk
drives should have provided with your equipment, or if that is not available,
use the CP/M PIP program.
Label the disks as backup disks and use them over and over.
It is always wise to keep your files backed up at least
one generation.
That way, should any problems arise the recovery procedure is quick
and relatively easy. (Be sure to make copies of your
backup files before you use them for
recovery.)  Do not minimize the importance of this step; adequate backups
are crucial to successful error recovery.  The information on your disks
represents a substantial investment in time and effort.
.sp2
.ce
SORT THE PAY TRANSACTION BATCH FILE
.pp
The Pay Transaction file is the only Payroll System file that must
be sorted for payroll application
processing.  To sort it you simply choose
the sort program from the menu (or if your CP/M compatible operating system
does not have the SUBMIT command, you run the sort outside the Payroll
System; see Chapter 26).
You must have room for the sorted output and temporary workfiles
on the same disk as the input file.  The QSORT sorting program requires
a sort parameter file which the Payroll System creates
automatically at first time startup
after the Parameter Entry, Account Entry, and System Deduction Entry steps
are complete.  Whenever you change the
location of the Pay Transaction file through the Parameter Entry program,
you must run
the Create Sort Parameters program to generate a new sort parameter file.
The Create Sort Parameter File program issues no requests.  Be sure to print
the audit proof before you sort the batch.
.sp2
.cp6
.ce
RUN THE APPLY PAYROLL PROGRAM
.pp
The Apply Payroll program
separates transactions for employees to be paid this application from
those to be paid later, creates a check file with the
earnings and deduction
information needed for each check, updates the company history and
employee history files, and creates a General Ledger
batch file (if GL interface
is requested).	The Apply Payroll program consumes all the
records in the Pay Transaction file and creates a new Pay Detail file that
contains detailed transaction information both for employees
paid on this application (so the Payroll Journal can be printed)
and employees to be paid on the next application.
.sp2
.ce
PRINT THE PAYROLL JOURNAL
.pp
Before you can run the Apply Payroll program again, you must print
at least one copy of the Payroll Journal.  This report is part of the
official audit trail. It is a detailed record of the pay transactions
for the employees paid this application. It also summarizes company-wide
earning and deduction information, and presents a ledger account
summary.
.sp2
.ce
PRINT PAYCHECKS AND/OR THE PAYCHECK REGISTER
.pp
After you complete the previous steps, you print the paychecks
and the Paycheck Register.  If you have not requested
computer check printing,
print the check register and take
the information you need from it to prepare the paychecks manually.
The paychecks produced by the Payroll System check
writing program are designed to use Moore Business Form paychecks.
.he
.bp
.ig
CHAPTER 13
E N T E R I N G    P A Y    T R A N S A C T I O N S
.sp2
.pp5
The Pay Transaction Entry program is used to enter employee time,
commissions, tips, bonuses, sick pay, vacation pay, advances, and
other types of pay.  It is also used to enter vacation, sick leave,
and comp time taken, reimbursement for expenses, repayment of advances,
and additional withholding when requested.  Detailed instructions on
how to enter the various transactions
can be found in Chapter 26.  This chapter presents a general
overview of the Pay Transaction Entry program's facilities.
.he Reference: Chapter 13
.pp
You may enter any number of transactions, in any order, for any active
employee.  For example, if some of your employees are paid by the
hour once a week on Tuesday for the week before,
and some are paid a salary once
a month on the first Tuesday of the following month, transactions
may be entered into the same batch for both groups.
.sp
.cp9
.li

		 |	month 1 	 |   month 2
		 |week1|week2|week3|week4|week1|week2|
		 |mtwtf|mtwtf|mtwtf|mtwtf|mtwtf|mtwtf| etc...
		 -------------------------------------
PAYDAY (weekly) 	 t     t     t	   t	 t
PAYDAY (monthly)			   t

	Figure 13.1: Weekly employees paid every Tuesday.
		     Monthly employees paid first Tuesday.
.sp
The hours worked by the hourly employees can be entered at the end of each
day, or at the end of the week,
as long as you enter them some time before the payroll
application and paycheck printing scheduled for Tuesday.  Special
transactions for hourly employees, such as
reimbursement for expenses, comp time earned,
advances, etc., can be entered at any time until the Tuesday
deadline.  Transactions for monthly employees can be entered any time
during or after the month, as long as you  enter them before the
first Tuesday of the next month, which is when the payroll application
program is scheduled to be run and paychecks printed.
.pp
Since salaried
employees normally receive a regular amount each pay period, a Rate1
pay transaction is usually not required, although special transactions
may be needed.	When no Rate1 transaction has been entered for an
active employee, the Apply Payroll program automatically generates
a transaction for the Rate1 amount times the number of Autogen
Units established for the employee
(usually 1 unit for monthly employees, 40 for an hourly employee paid
weekly, etc.).
.pp
In the example above, each
Tuesday the Apply Payroll program is run once for the weekly employees.
On the first Tuesday of the second month you would
run it twice, once for the week-based
employees, and once for the month-based employees.  The paycheck,
check register, and payroll journal print programs are also run
once for week-based and once for month-based employees.
.pp
The Pay Transaction Entry program requests four items of information
for each transaction: 1). the number of the employee who
is affected, 2). the transaction type or
code, 3).
the number of units (in hours, days, dollars, pieces, etc.), and 4).
a date, either the date the work was
done, vacation taken, bonus given, etc., or the date the transaction is
entered onto the file.	The date you give is your decision, although
if none is entered, the run date given at startup is used in default.
This date is for reference only.  The date used as the General Ledger
transaction reference is the check date entered when you begin the
Apply Payroll program.
.pp
Below is a complete list of the transactions you may enter through
the Pay Transaction Entry program.  Beside each is the transaction (or
rate) code you would enter to indicate the transaction type.
.in5
.sp
.ti-5
(0)  AUTOMATIC CALCULATION OF REGULAR AND OVERTIME PAY:  When you  enter
a Rate0 transaction and the number of units,
the Payroll System automatically calculates pay at Rate2
for units over the normal number of units, and Rate1 for the normal units.
For example, if Normal Units is set to 40 for an hourly employee, a Rate0
transaction of 48 units causes pay for the first 40 units to be calculated
at Rate1, and pay for the next 8 units to be calculated at Rate2.  This
transaction type is intended as a convenience. The same thing could be
accomplished by two transactions, a Rate1 transaction for 40 hours, and
a Rate2 transaction for 8 hours.
.sp
.ti-5
(1)  RATE1 PAY: When you enter a Rate1 transaction, the employee
is paid at Rate1 for the given number of units.  Rate1 is normally
Regular Pay, but can be given any description through the Parameter
Entry program.	Entering a Rate1 transaction always
overrides (suppresses) automatically
generated pay. If a Rate1 transaction is not entered, the Rate1 pay
amount is calulated according to the number of Autogen Units set
for the employee through the Employee Rate Entry program.
Rate1 pay is subject to all taxes and percentage-based
deductions that apply to the
employee.
.sp
.ti-5
(2)  RATE2 PAY: When you enter a Rate2 transaction, the employee
is paid at Rate2 for the number of units entered.  Rate2 is normally
overtime pay, but can be given any description through the Parameter
Entry program. Rate2 pay is subject to all taxes and percentage-based
deductions that
apply to the employee.
.sp
.ti-5
(3)  RATE3 PAY: When you enter a Rate3 transaction, the employee
is paid at Rate3 for the number of units entered.  Rate3 is normally
special overtime pay, but can be given any description through the
Parameter Entry program. Rate3 pay is subject to all taxes and
percentage-based deductions that apply to that employee.
.sp
.ti-5
(4)  RATE4 PAY: When you enter a Rate4 transaction, the employee
is paid at Rate4 for the number of units entered.  Rate4 is normally
used to enter commission pay, piecework pay, or "other" pay,
but can be given any
description through the Parameter Entry program. If a
laborer who receives 65 cents per box (Rate4) for picking fruit
picks 50 boxes,
you would answer 4 at the
request for the transaction code, and 50 for the number of units.
If a salesperson who is paid once a month for 1 unit at Rate1
(for example, 1 unit at $1,500 per unit) is to receive in addition to his
or her regular pay a commission of $500
for closing a large sale, you would set Rate4 for the employee to
one dollar ($1.00) and enter a Rate4 transaction at 500 units, in
addition to entering the normal Rate1 transaction.
Rate4 pay is taxed according to the exclusion code set for Rate4 pay
in Parameter Entry.
.sp
.ti-5
(5)  VACATION TAKEN: To record the fact that an employee has used
some of his or her accumulated vacation time, you would enter
a transaction with a transaction code of 5. The number of units
would be the number of days or hours the employee took for vacation.
.sp
.ti-5
(6)  SICK LEAVE TAKEN:	To record the fact that an employee has used
some of his or her accumulated sick leave, you would enter a
transaction for that employee with the transaction code of 6.  The
number of units would be the number of days or hours the employee
took in sick leave.
.sp
.ti-5
(7)  COMP TIME TAKEN: To record the fact that an employee has used
some of his or her accumulated comp time, you would enter a
transaction for that employee with the transaction code of 7.  The
number of units would be the number of days or hours the employee
took in compensatory time off.
.sp
.ti-5
(8)  COMP TIME EARNED: To record the fact that an employee has
earned compensatory time off, you would enter a transaction for
that employee with the transaction code of 8.  The number of units
would be the number of hours or days the employee has earned.
.sp
.ti-5
(9)  BONUS: To give an employee a bonus, you would would
enter a transaction for that employee with the transaction code of
9.  The number of units would be the dollar amount of the bonus.
If all or most employees are to receive the bonus, if the bonus is to be
distributed over several pay periods, or if the amount is a
percentage of an employee's salary, you will probably want to use the
System Deduction Entry program or the Employee Deduction/Earnings and
Cost program to order the bonus.  Bonus transactions entered through
the Pay Transaction Entry program are subject to all taxes and
percentage-based deductions
that apply to the employee.
.sp
.ti-5
(10) TIPS COLLECTED: To pay an employee for tips collected by the employer
you would enter a transaction with the transaction code of 10. The
number of units would be the dollar amount of the tips owed by the employer
to the employee.  Tips are added to the employee's other earnings, and are
subject to all taxes and percentage-based
deductions that apply to the employee, including
FICA.  Since the employer is not required to pay FICA tax on employee
tips, tips are excluded from the employee's earnings for the purpose of
calculating the employer's FICA tax contribution.
.sp
.ti-5
(11) ADVANCE: To pay an employee an advance on his or her salary, you would
enter a transaction with a transaction code of 11.  The number of units
would be the dollar amount of the advance.  An advance is not subject
to any taxes or percentage-based deductions. It is distributed to the
cost accounts set for the employee.  An adjusting General Ledger entry may
be required.
.sp
.ti-5
(12) SICK PAY: To mark employee pay as sick pay, you would enter a transaction
with a transaction code of 12.	The number of units would be the dollar
amount of the sick pay.  Sick pay is subject to all taxes and
percentage-based
deductions that apply to the employee.
.sp
.ti-5
(13) VACATION PAY: To mark employee pay as vacation pay, you would enter
a transaction with a transaction code of 13.  The number of units
would be the dollar amount of the vacation pay.  Vacation pay is
subject to all taxes and percentage-based deductions that apply to the
employee.
.sp
.ti-5
(14) OTHER EXCLUDED PAY: Pay transactions with a transaction code of 14
are excluded from all taxes and percentage-based
deductions that apply to the employee.
.sp
.ti-5
(15) EXPENSE REIMBURSEMENT: To reimburse an employee for expenses, you
would enter a transaction with a transaction code of 15.  The number
of units would be the dollar amount of the reimbursement.  The amount
is added to the employee's net earnings and is not subject to any taxes
or percentage-based
deductions that apply to the employee. It is, however, distributed to the
ledger accounts set for the employee.  An adjusting General Ledger entry
may be required.
.sp
.ti-5
(17) OTHER PAY: Pay transactions with a transaction code of 17 are subject
to all taxes and percentage-based deductions that apply to the employee.
.sp
.ti-5
(18) TIPS REPORTED: Tips reported are tips pocketed by the employee
when earned and later
reported to the employer.  Because Tips Reported are tips that have already
been collected by the employee, the employer does not pay the employee
for them.  However, these tips must be added to the employee's earnings
for tax purposes.  Tips Reported are included in the employee's gross taxable
earnings, but are also deducted from earnings since
the employee has already collected them.  The employer is exempt from
paying FICA tax on employee tips. Tips reported are not distributed to the
cost accounts set for an employee.
.sp
.ti-5
(27) ADVANCE REPAY: To deduct money from an employee's net pay for
the repayment of an advance on his or her salary, you would enter a
transaction with a transaction code of 27.  The number of units would
be the dollar amount of the repayment.
.sp
.ti-5
(28) FWT ADD-ON: If an employee requests additional Federal income tax
withholding,
you would enter a transaction with a transaction code of 28.  The number
of units would be the dollar amount of the additional withholding.
.sp
.ti-5
(29) SWT ADD-ON: If an employee requests additional state income tax
withholding,
you would enter a transaction with a transaction code of 29.  The number
of units would be the dollar amount of the additional withholding.
.sp
.ti-5
(30) LWT ADD-ON: If an employee requests additional local income tax
withholding,
you would enter a transaction with a transaction code of 30.  The number
of units would be the dollar amount of the additional withholding.
.sp
.ti-5
(31) FICA ADD-ON: If an employee requests additional FICA tax withholding,
you would enter a transaction with a transaction code of 31.  The number
of units would be the dollar amount of the additional withholding.
.qi
.he
.bp
.ig
CHAPTER 14
P A Y R O L L	 A P P L I C A T I O N
.he Reference: Chapter 14
.sp2
.pp5
The Apply Payroll program is the main processing program of
the Payroll System.  It calculates employee gross earnings by
multiplying the number of pay units
times the rate information stored on
the Employee Master file, and adding bonuses, tips, and other pay
transactions.  If no Rate1 units have been entered for an Active
employee, the program generates pay for that employee automatically,
according to the number of Autogen Units set for the employee (unless
set to zero units).  The Payroll Application program calculates income
tax withholding, and deductions for FICA, SDI, and
user-defined deductions.  The program updates registers of
historical information, such as quarter and year to date withholding
for employees, company liabilities, vacation used, etc.  It
creates a Check File that contains the information needed to produce
checks and print the Paycheck Register.  If the option has been selected,
the Apply Payroll program creates a General Ledger batch file for
entry into your General Ledger.  Finally, as a matter of housekeeping,
the program increase the apply number for the interval base (i.e.,
if the application was for week-based employees, it increases its counter
of the number of times the application program has been run for week-based
employees by one), and resets the Last Day of Last Period for the
intervals (pay periods) processed by the application.
.pp
Before you run the Payroll Application make certain that:
.sp
.in5
.ti-3
1. Transactions for employees to be paid by this run of
the Apply Payroll program have been entered into the
Pay Transaction batch file, and checked carefully
for accuracy and completeness.
If you forget someone, you must re-run the application using your
backup files since there is no provision for processing individual
employees separately.
You cannot just run another application since this would
throw off the Payroll System's internal count of the applications run
and distort the historical statistics accumulated so far.
.sp
.ti-3
2. You have backup copies of all of your processing disks.  Do not
neglect this step.
.sp
.ti-3
3. A final, audit proof of the Pay Transaction batch has been printed
to protect the audit trail.
If this has not been done, the application program will not let you
continue.
.sp
.ti-3
4. The Pay Transaction file has been sorted into employee number
order.
.qi
.pp
The Payroll System allows you to pay employees on any combination of
four pay schedules: weekly, bi-weekly, monthly, and semi-monthly.
The weekly and bi-weekly schedules (or "intervals") are called week-based
intervals, and the semi-monthly and monthly schedules are called
month-based intervals.	Week-based employees and month-based employees
are never processed by the same run of the Apply Payroll program,
although transactions for both kinds of employees can exist on the
Pay Transaction file.  The weekly and semi-monthly pay intervals are
referred to as the "fast" intervals because payroll application occurs
twice as often for them as for the so-called "slow"intervals, bi-weekly
and monthly (there may be 52 or 53 weekly applications in one year, and
only 26 or 27 bi-weekly applications; payroll application occurs 24 times
a year for semi-monthly employees, while application takes place only
12 time a year for the "slow" monthly employees).
.pp
The Payroll System maintains two counters of the number of
applications that have been run.  One counter keeps track of week-based
applications, the other of month-based applications.  When "fast"
interval employees are on the system, each run of the application program
increases the counter for that interval base by one.  When "slow" interval
employees are on the system, each run of the application program increases
the counter for that interval base by two.  When both fast and slow intervals
of the interval base are used (i.e., semi-monthly and monthly, weekly and
bi-weekly), every other application processes both fast and slow employees,
which should occur only when the application number is even. When you add
new intervals to the system, they should be timed to coincide with the
the proper odd or even number.
.pp
In general, only fast interval employees are processed on odd applications.
An even application processes both fast and slow interval employees (or
just slow interval employees if the fast inteval is not used).	The
payroll application for month-based employees
that takes place at the end of the month for
work done during the month (the whole month for
monthly employees, the last half for semi-monthly) should always be even.
The application number for week-based applications will, of course,
be even every two weeks.
.pp
When you run the Payroll Application program, before beginning it asks
you for a Check Date.  This is the date that will print on the checks,
and thus the date you legally incur liabilities for such things as
Federal, state, local and FICA withholding.  This is also the date
on the General Ledger transactions.  The Apply Payroll program
checks this date carefully to make sure end of period processing
is not required. If it is, you are not permitted to run the application
until you have completed end of period processing.
.he
.bp
.ig
CHAPTER 15
P R I N T I N G    P A Y C H E C K S
A N D	 T H E	  P A Y C H E C K    R E G I S T E R
.sp2
.pp5
The Payroll System paychecks are designed to print on Moore Business
Form payroll checks; contact your local representative.
Earning and Deduction information
prints in 16 boxes on the check stub. Box 1 holds the Gross Pay amount,
boxes 2 through 7 hold other earnings categories, box 8 contains the
Net Pay amount, and boxes 9 through 16 are reserved for deductions.
A complete list of which values are assigned to which boxes appears
below in Figure 15.1.
.he Reference: Chapter 15
.pp
When you create a employee specific or systemwide earning or deduction
category, you are asked to enter the check category where the amount
should be printed. Boxes (check categories) 2 through 7 are set aside
for earnings, while boxes 9 through 16 are set aside for deductions.
You cannot assign a deduction to an earnings box or vice versa.
When more than one item is assigned to a box, they are added together.
.pp
Earning or deduction amounts generated by payroll
transactions are assigned to stub boxes automatically by the system
and cannot be changed, though that does not prevent you from assigning
other items to those boxes.
.sp
.cp30
.li
	  STUB BOX		 ITEM

 EARNINGS    1		     Gross Earnings
	     2		     Rate1 Pay (Regular)
	     3		     Rate2 Pay (Overtime)
	     4		     Rate3 Pay (Special Overtime)
	     5		     Rate4 Pay (Commission)
	     6		     Other Pay transaction
			     Other Excluded Pay transaction
			     Bonus transaction
			     Vacation Pay transaction
			     Sick Pay transaction
			     EIC Payment
			     Expense Reimbursement transaction
	     7		     Tips Collected
	     8		     Net Pay

 DEDUCTIONS  9		     Federal Withholding
			     FWT Add On transaction
	    10		     State Withholding
			     SWT Add On transaction
	    11		     Local Withholding
			     LWT Add On transaction
	    12		     FICA Withholding
			     FICA Add On transaction
	    13		     SDI Withholding
	    14		     Tips Reported transaction
			     Advance Repayment
	    15		     [no system assigned significance]
	    16		     [no system assigned significance]
.pp
To print the checks after the Apply Payroll program has been run and
a Check file created, you call the Print Pay Checks program from the menu.
It will tell you what the number of the last check printed was, and
ask you to enter the Starting Check Number for this run (it can be
any number you desire).  The program lets you print a test form
to check your form alignment.
.he
.bp
.ig
CHAPTER 16
E N D	O F   P E R I O D   P R O C E S S I N G
.sp2
.pp5
At the end of each quarter you must print the Federal 941 Report,
the Employer's Quarterly Federal Tax Return and run the Close
for the Quarter program.  At the end of each year you must print
the Federal 940 Report,  W-2 forms
for all employees, and run the Close for the Year program.
.he Reference: Chapter 16
.sp2
.ce
ENDING THE QUARTER
.pp
The quarter end processing steps must be undertaken sometime after
the last application of the last quarter, and before the first payroll
application of the next quarter.  The Apply Payroll program does not
let you continue if the Check Date you enter is past the last day
of the quarter (Figure 16.1), and the Print 941
Report and Close for Quarter programs
have not been run.  You are  also prevented from continuing if not enough
applications have been run for the quarter to end (3 monthly, 6 semi-monthly,
at least 13 weekly, and at least 6 bi-weekly).	New intervals may only
be added after quarter end processing and before the first application of
the next quarter.
.sp
.cp6
.li
	      QUARTER		LAST DAY OF QUARTER

		 1		    March 31
		 2		    June 30
		 3		    September 30
		 4		    December 31

		      (Figure 16.1)
.pp
The 941 Report may be printed directly onto the government form, or
you can print it on plain paper and transfer the information by hand.
You can suppress the automatic printing of your company name and address
if you wish to use the printed label provided by the government.
The Print 941 Report program displays a facsimilie reproduction of the
941 report with the current figures, and allows you to change several
of the items displayed.  You can enter an adjustment of withheld income
tax from previous quarters, an adjustment of FICA taxes, or change
the payment dates and amounts if necessary to correspond to the actual
dates and payment amounts. The report may be printed as many times as
you desire.
.pp
The Close for Quarter program resets various registers of historical
information to zero in preparation for beginning the next quarter.  This
program issues no requests.  You must print the 941 report before you are
permitted to run the Close for the Quarter program.
.sp
.ce
ENDING THE YEAR
.pp
The year end processing steps must be undertaken sometime after the Close
for the Quarter program has been run for the fourth quarter, and before
the first run of the Apply Payroll program for the new year.
The Apply Payroll program does not let you continue if the Check
Date you enter is past the first of the year and the Close for the Year
program has not been run.  The Close for the Year program may not be
run until the Federal 940 Report and the W-2 forms have been printed.
And the Federal 940 Report and the W-2 forms may not be printed until
the Close for the Quarter program has been run for the fourth quarter.
.sp
.li
		   YEAR END PROCESSING STEPS

		1. Close for the fourth quarter.
		2. Print Federal 940 Report.
		3. Print W-2 forms for all employees.
		4. Run Close for the Year program.

		 (Figure 16.2: Ending the Year)
.pp
Print the 940 report on plain paper and transfer the information to
the government form.  There are no requests issued by this program.
.pp
The W-2 forms print in the standard W-2 format.  The Print
W-2 Report program lets you specify the number of forms to print for each
employee.  This option makes it possible to print W-2's  on multiple copy
carbon forms, or print the required number of W-2 forms on carbonless forms.
You can suppress printing of local and state withholding amounts for
the run, or print Federal, state, and local withholding information together.
The program also lets you set the starting print column if you wish.
.pp
You can produce W-2 forms for all employees, a range of employees,
terminated employees only, or for a single employee.  The Print W-2
Report program may be run as many times as needed before you run the
Close for the Year program.
.pp
The Close for the Year program resets various registers of
historical information in preparation for the new year.  It issues
no requests.  It does not run until the 940 Report has been printed
and W-2 forms have been printed for all employees.
.he
.bp
.ig
CHAPTER 17
O P T I O N A L    R E P O R T S
.sp2
.pp5
The reports below (Figure 17.1)
are optional in that their production is not
a prerequisite to further processing steps.  They can be produced
at any time.
.he Reference: Chapter 17
.sp
.li
	       1. Account Print
	       2. Employee Master List
	       3. Employee History Report
	       4. Vacation Report
.pp
The Account Print is a short report that lists the currently
valid Payroll System accounts.	To produce the Account Print you
align the paper in your printer, and
choose the Print Accounts program from the menu.
.pp
The Employee Master List can be produced in two versions.  One is
a short version that gives an employee's name, number, social
security number, pay interval, status (active, terminated, on leave),
pay type (hourly, salaried), phone number, start date, birth date,
Rate1 rate, and job code.  It prints one line for each employee.
.pp
The full version of the Employee Master List devotes one page
to each employee.  It lists all of the above, plus information
on pay rates, vacation, sick leave, comp time, withholding allowances,
tax exemptions, employee specific earnings and deductions, and
cost distribution.
.pp
When you select the Print Employee Master List program from the system menu,
another menu appears from which you can select the desired printing
option.  This menu is reproduced below (Figure 17.2).
.sp
.cp19
.li
		 Y O U R    C O M P A N Y    N A M E
			 AC PAYROLL SYSTEM
00/00/00	   EMPLOYEE MASTER LIST PRINT MENU

    1. ALL EMPLOYEES			8. ALL HOURLY EMPLOYEES
    2. ALL WEEKLY EMPLOYEES		9. ALL SALARIED EMPLOYEES
    3. ALL BI-WEEKLY EMPLOYEES	       10. ALL ACTIVE EMPLOYEES
    4. ALL SEMI-MONTHLY EMPLOYEES      11. ALL ON-LEAVE EMPLOYEES
    5. ALL MONTHLY EMPLOYEES	       12. ALL TERMINATED EMPLOYEES
    6. ALL WEEK BASED EMPLOYEES        13. ALL DELETED EMPLOYEES
    7. ALL MONTH BASED EMPLOYEES       14. A RANGE OF EMPLOYEES

	       ENTER YOUR SELECTION NUMBER: [  ]

	       RANGE FROM [	] TO [	   ]

	       FULL INFORMATION OR SINGLE LINE (F OR S) [ ]

			  (Figure 17.1)
.pp
The Employee History Report presents summarized quarter to date and
year to date earnings and deduction information.  To print this report
you choose the Print History Report from the system menu.
.pp
The Vacation Report includes information on sick pay and comp time.
To print this report you select it from the system menu.
.he
.bp
.ig
CHAPTER 18
I N T E R A C T I N G	 W I T H    T H E    S Y S T E M
.sp2
.pp5
This chapter explains how to interact with
the Payroll System, and provides
additional information that may be useful to
those new to microcomputer operation.
.he Operations: Chapter 18
.sp2
.ce
HOW TO ANSWER A REQUEST
.pp
First, move the cursor to the bracketed field that corresponds to
the request.  (The cursor is the moving spot of light on the screen,
usually triangular or rectangular in shape.  It may or may not be
blinking, depending on your equipment.)  If the field is empty,
type in your response, and then press the RETURN key.
If the field is full, type in your response over the current value
and press RETURN.  For example,
.sp
.cp5
.li
    COMPANY NAME: [Universal Enterprises       ]  <Current Value>

    COMPANY NAME: [Universal Inc.rprises       ]  <Your Entry>

    COMPANY NAME: [Universal Inc.	       ]  <After RETURN>
.sp
If the current value is acceptable, just press RETURN.
.sp2
.ce
HOW TO MOVE FROM REQUEST TO REQUEST
.pp
After you type in your response to a request and hit RETURN, the
cursor moves along a fixed path to the next request, or around
to the first request.  To change
the value currently in a field, you press RETURN to move the cursor
forward to the desired request, and type in your response.  To move
back one field (that is, one request) in a Display Management
System entry screen,
give the command:
.sp
.li
		CTRL-B
.sp
Note, however, that when the cursor is at a field labeled Employee Number
or Record number, the CTRL-B command means something different.  At those
locations CTRL-B causes the program to display the previous record on the
file (i.e., "go BACK and get the previous record").
.sp2
.cp6
.ce
HOW TO "THUMB" THROUGH A FILE
.pp
When you select the Employee Entry, Employee Ded/Earn and Cost, or the
Employee Rate Entry programs, the appropriate entry screen appears with
the cursor waiting at the Employee Number request.  If there are records
on the Employee Master file, you can "thumb" through the file by any
combination of the following three methods:
.sp
.in3
.ti-3
1. Give the command CTRL-N to get the NEXT record on the file.	When
no records have been accessed, this command gets the first record on the
file.
.sp
.ti-3
2. Give the command CTRL-B to go BACK and get the previous record on the
file.  When no records have been accessed, this command brings up the last
record on the file.
.sp
.ti-3
3. Enter the five digit employee number to get the record for that employee.
.qi
.pp
When you select the Pay Transaction Entry, or the Account Entry programs
the appropriate entry screen appears with the cursor waiting at the
Record Number request.	If valid records exist on the file you can
"thumb" through the file by any combination of the following three
methods:
.sp
.in3
.ti-3
1. Give the command CTRL-N to get the NEXT record on the file.	When
no records have been accessed, this command gets the first record on
the file.
.sp
.ti-3
2. Give the command CTRL-B to go BACK and get the previous record on
the file.  When no records have been accessed, this command gets the
last record on the file.
.sp
.ti-3
3. Enter the "relative record number" of the record you want to access.
The relative record number is the position of the record relative to the
beginning of the file.	For example, the first record on the file has
relative record number 1, the second has relative record number 2, etc.
The Payroll System Account Reference Number is the same thing as the
relative record number of an Account File record.
.qi
.sp2
.cp6
.ce
HOW TO FIX A TYPOGRAHPICAL ERROR
.pp
To back up and erase one character before pressing RETURN, give
the command:
.sp
.li
	     CTRL-H
.sp
as many times as necessary to back up to the erroneous entry, then
type in the correct characters.  Hit RETURN to continue.
.pp
If you've already pressed RETURN, give
the command CTRL-B to go back to the previous field and type in the
correct response.
.sp2
.cp6
.ce
HOW TO CANCEL AN ENTRY
.pp
While adding an employee record through the Employee Entry program,
an account through the Account Entry program, or a pay transaction
through the Pay Transaction Entry program, before hitting ESCAPE to
begin the next record, you can cancel the entire employee, account,
or transaction record by giving the command:
.sp
.li
		CTRL-X
.sp
Once the entry has been completed and entered onto the system by hitting
ESCAPE, the cancel command has no effect.
.sp2
.cp4
.ce
HOW TO GIVE A CONTROL COMMAND
.pp
To give a CONTROL command, such as CTRL-H or CTRL-N, you hold down
the CONTROL key (CTRL or CTL on some keyboards) as you would a SHIFT
key on a typewriter, and while holding it down, type the appropriate
letter (either upper or lower case).  You can depress the CONTROL key
for as long as it takes to enter the command.
.sp2
.cp6
.ce
HOW TO ENTER A NUMBER
.pp
When you answer a request with a mult-digit number or dollar amount,
do not type in the commas, if any.  The program inserts them where needed.
Two is the largest number of digits allowed to the right of the decimal
point.	You must type in the decimal point (period) when you enter a
decimal number.  When you enter a flat dollar amount (no cents), you need
not type in the period and two zeros, the program inserts them for you.
.sp
.cp8
.li
     EXAMPLE: To enter $4,520.00 into the MAX PAY field,
	      type in:

		  MAX PAY: [4520	]

	      After hitting RETURN, the field shows,

		  MAX PAY: [	4,520.00]
.sp
.cp8
.li
     EXAMPLE: To enter $4,520.35 into the MAX PAY field,
	      type in:

		  MAX PAY: [4520.35	]

	      After hitting RETURN, the field shows,

		  MAX PAY: [	4,520.35]
.sp2
.cp6
.ce
HOW TO ENTER A DATE
.pp
The Payroll System checks all date entries to make sure they
are valid dates and conform to the date format (MM/DD/YY or DD/MM/YY)
specified in
Parameter Entry.  The system bars you from entering
a date such as April 31 (since April has only 30 days).
Nor will it let you enter February 29 for a non-leap-year.
The American date format (MM/DD/YY) is used unless you specifically
request the international format (DD/MM/YY).  You need not type leading
zeros for single digit months or days, but you are required to type
in the slash marks (/).
.sp
.li
      EXAMPLE: To enter March 5, 1980, you would type:

		   [3/5/79  ]

	       and then RETURN.  After RETURN, the field
	       would show:

		   [03/05/79]
.sp2
.cp6
.ce
HOW TO HANDLE THE "FLOPPY" DISKS
.pp
NEVER touch the magnetic material exposed through the protective sleeve.
.pp
NEVER let the disk come within the force of a magnetic field (such as
the tip of a magnetized screwdriver).
.pp
NEVER remove a disk from a disk drive while the "access" light is on.
Wait until the CP/M prompt character is showing, or the program instructs
you to remove it.
.pp
NEVER turn the power on or off while a disk is inserted in the machine.
.pp
ALWAYS insert the disk into the drive with the exposed "slot" pointing
toward the rear of the machine.  The direction in which the label should
face depends on your equipement.
.pp
ALWAYS put the disk back in its protective sleeve and store it safely
away when you are through using it.
.sp2
.cp6
.ce
HOW TO END THE CURRENT FUNCTION OR PROGRAM
.pp
Hitting the ESCAPE key takes you one "layer" out of the system.  For
example, while adding employees to the Employee Master file in the
Employee Entry program, ESCAPE signals that you have completed the record
for that employee and are ready to add another.  An ESCAPE right after
the previous ESCAPE (before you begin the next employee
record), takes you out of the Add-Records mode, and returns the
cursor to the Employee Number field where you can give other commands
(such as CTRL-N or CTRL-B).  A third ESCAPE ends the Employee Entry
program and returns you to the main system menu.  ESCAPE at the system
menu ends the Payroll System altogether.  When the CP/M prompt appears
and the access light goes out, you can remove the disks.
.sp2
.cp3
.ce
HOW TO END A PRINTOUT
.pp
If your paper should jam in the printer, or if you want to end a
printout for some other reason, just press any key on the keyboard.
To restart where you left off, type "RUN".  To end the program, type
"STOP".

